\begin{table}[t]
\centering
\captionsetup{justification=raggedright,singlelinecheck=false}
\caption{
\textbf{Differential-attrition tests, by study.} \\ 
Per-study chi-square test of attrition status against assigned condition on the matched-only sample. Cells are the per-arm rates reported in the preceding attrition table; Constrained AI is treated as a distinct arm in Study 2.
} 
\begin{tabular}{lccccc}
\toprule
\multicolumn{1}{c}{\textbf{Study}} & \multicolumn{1}{c}{\textbf{k arms}} & \multicolumn{1}{c}{\textbf{chi-square}} & \multicolumn{1}{c}{\textbf{df}} & \multicolumn{1}{c}{\textbf{p}} & \multicolumn{1}{c}{\textbf{n matched}} \\ 
\midrule\addlinespace[2.5pt]
\noalign{{\color{gray!15}\hrule height \dimexpr\ht\strutbox+\dp\strutbox\relax depth 0pt}\nointerlineskip\vskip -\dimexpr\ht\strutbox+\dp\strutbox\relax}
Study 1 & 5 & 11.35 & 4 & 0.023 & 9,776 \\ 
Study 2 & 4 & 17.51 & 3 & <.001 & 4,553 \\ 
\noalign{{\color{gray!15}\hrule height \dimexpr\ht\strutbox+\dp\strutbox\relax depth 0pt}\nointerlineskip\vskip -\dimexpr\ht\strutbox+\dp\strutbox\relax}
Study 3 & 3 & 0.03 & 2 & 0.986 & 2,569 \\ 
Study 4 & 3 & 1.29 & 2 & 0.525 & 2,995 \\ 
\bottomrule
\end{tabular}
\label{tab:attrition_chisq}
\end{table}
